\documentclass[UTF8,11pt]{ctexart}
\usepackage{amsmath, amssymb, amsfonts}
\usepackage{graphicx, geometry, xcolor}
\usepackage{algorithm, algorithmic}
\usepackage{bm}
\usepackage{hyperref}
\usepackage{titlesec}
\usepackage{enumitem}
\usepackage{tabularx}
\usepackage{array}
\usepackage{fancyhdr}

\geometry{a4paper, margin=0.8in}

% Set global no indentation
\setlength{\parindent}{0pt}

% Set hyperlink style
\hypersetup{
    colorlinks=true,
    linkcolor=blue,
    urlcolor=blue,
    citecolor=blue
}

% Set header and footer
\pagestyle{fancy}
\fancyhf{}
% \fancyfoot[C]{\thepage}  % Comment this line to remove page numbers
\renewcommand{\headrulewidth}{0pt}
\renewcommand{\footrulewidth}{0pt}  % Add this line to remove footer line

% Custom title format
\titleformat{\section}{\Large\bfseries}{\thesection}{1em}{}[\titlerule]
\titlespacing*{\section}{0pt}{10pt}{7pt}

% Define colors
\definecolor{darkblue}{RGB}{0,51,102}

\begin{document}

%================= Header =================
\begin{center}
    {\LARGE \textbf{Guanting Liu}}\\[0.4em]
    \href{liuguanting@gdiist.cn}{liuguanting@gdiist.cn} $|$ +86 13302629998 $|$ \href{https://scholar.google.com/citations?hl=en&user=VUuphxEAAAAJ&view_op=list_works&sortby=pubdate}{Google Scholar}
\end{center}

\vspace{1.0em}

%================= Bio =================
\section*{Bio}
\textbf{The Hong Kong Polytechnic University} \hfill 2026-Now\\
\text{PhD Student}

\vspace{0.3em}
\textbf{Guangdong Institute of Intelligence Science and Technology} \hfill 2025-2026 \\
\text{Research Assistant}

\vspace{0.3em}
\textbf{Nanjing University of Information Science and Technology} \hfill 2021-2025 \\
\text{Bachelor of Science}




\section*{Research Interests}
Computational neuroscience, cognitive neuroscience, brain-computer interface, deep learning.




\section*{Main Papers}
\hfill $^1$ indicates co-first author, $^*$ indicates corresponding author.

[1] \textbf{Guanting Liu$^1$}, Ying Yan$^1$, Jun Cai, Edmond Qi Wu, Shencun Fang$^*$, Adrian David Cheok, \& Aiguo Song. (2025). "GCD: Graph contrastive denoising module for GNNs in EEG classification". \textit{Expert Systems with Applications}, 265, 126013. \href{https://doi.org/10.1016/j.eswa.2024.126013}{[DOI]}

\vspace{0.5em}
[2] \textbf{Guanting Liu$^*$}, Ying Yan, Jun Cai, Adrian David Cheok, Edmond Qi Wu, \& Aiguo Song. "Rethinking Kalman Filters for Motor Brain-Machine Interface: The Fundamental Limitations and A Perspective Shift". \textit{biorxiv preprint}. \href{https://doi.org/10.1101/2025.07.10.664247}{[DOI]}

\vspace{0.5em}
[3] \textbf{Guanting Liu$^*$}, Ying Yan, Jun Cai, Adrian David Cheok, Edmond Qi Wu, \& Aiguo Song. "A More Rational and Efficient Kalman Filter Design for Motor Brain-Machine Interfaces". \textit{2025 47th Annual International Conference of the IEEE Engineering in Medicine \& Biology Society (EMBC)}. \href{10.1109/EMBC58623.2025.11251710}{[DOI]}

\vspace{0.5em}
[4] \textbf{Guanting Liu$^*$}, Ying Yan, Sizhen He, Jun Cai, Adrian David Cheok, Edmond Qi Wu, \& Aiguo Song. "A Neuromorphic Approach for Brain-Machine Interface Using Spiking Neural Networks". \textit{2025 47th Annual International Conference of the IEEE Engineering in Medicine \& Biology Society (EMBC)}. \href{10.1109/EMBC58623.2025.11254255}{[DOI]}


\section*{Other Papers}

[5] Ying Yan$^1$, \textbf{Guanting Liu}$^1$, Jun Cai, Na Liu, Shencun Fang$^*$, Adrian David Cheok, Edmond Qi Wu, Chengcheng Hua, Aiguo Song. (2026). "Residual multi-dimensional Taylor network for epileptic electroencephalography detection". \textit{Engineering Applications of Artificial Intelligence}.

\vspace{0.5em}
[6] Ying Yan$^1$, \textbf{Guanting Liu$^1$}, Haoyang Cai, Edmond Qi Wu$^*$, Jun Cai, Adrian David Cheok, Na Liu, Tao Li, \& Zhiyong Fan. (2024). "A review of graph theory-based diagnosis of neurological disorders based on EEG and MRI". \textit{Neurocomputing}.

\vspace{0.5em}
[7] Ying Yan$^1$, \textbf{Guanting Liu$^1$}, Jun Cai$^*$, Edmond Qi Wu, Adrian David Cheok, Aiguo Song. A Novel Interpretable Multi-Layer Voting Network for Fault Diagnosis. (2026).

\vspace{0.5em}
[8] Ying Yan$^*$, \textbf{Guanting Liu}, Jun Cai, Edmond Qi Wu, Adrian David Cheok, Qiming Sun, Ming Li, \& Ying Zhou. (2023). "A Novel ResMTN-Based Methodology for HVAC AHU Fault Diagnosis". \textit{2023 Global Reliability and Prognostics and Health Management Conference (PHM-Hangzhou)}.

\vspace{0.5em}
[9] Haoyang Cai, Ying Yan$^*$, \textbf{Guanting Liu}, Jun Cai, Adrian David Cheok, Na Liu, Chengcheng Hua, Jing Lian, Zhiyong Fan, Anqi Chen. (2024) "WKLD-based feature extraction for diagnosis of epilepsy based on EEG". \textit{IEEE Access}.

\vspace{0.5em}
[10] Ying Yan, Haoran Li, Jun Cai, \textbf{Guanting Liu}, Edmond Qi Wu, Hao Wang, Chengcheng Hua, Yaowen Yu$^*$, Aiguo Song. (2026) "A novel correlation-driven cross-term compression polynomial network for classifying motion sickness levels". \textit{Engineering Applications of Artificial Intelligence}.



\section*{Research \& Projects}

\textbf{Large-scale human connectome-based task-driven brain modelling} \hfill 2026-Now
\begin{itemize}[leftmargin=1.0em, topsep=0pt, itemsep=0pt]
    \item Constructed a mesoscale MRI connectome-based structural scaffold, modeled regional dynamics with firing-rate units, and trained the network to perform cognitive tasks.
    \item Investigated efficient signal propagation and reliable information routing across the whole brain.
\end{itemize}

\textbf{Multi-levels cognitive computational modelling} \hfill 2025-Now
\begin{itemize}[leftmargin=1.0em, topsep=0pt, itemsep=0pt]
    \item Modeled biologically constrained models like leaky integrate-and-fire, Hodgkin-Huxley and analyzed their dynamics in cognitive tasks such as working memory and decision making.
    \item Utilized dynamics tools such as neural manifolds to analyze and compare the computational and neuroscientific principles of networks across different levels.
\end{itemize}

\textbf{Whole-brain simulation modelling of drosophila and zebrafish} \hfill 2025-2026
\begin{itemize}[leftmargin=1.0em, topsep=0pt, itemsep=0pt]
    \item Constructed a whole-brain model of the drosophila by incorporating its real connectome at a scale of 140,000 neurons, achieving predictions of resting-state neural activity.
    \item Modeled and predicted neural activity in zebrafish by incorporating single-neuron-resolution calcium imaging data, and reconstructed and analyzed its sensory-motor functions.
\end{itemize}

\textbf{Invasive motor brain-computer interface decoding algorithms designing} \hfill 2024-2025
\begin{itemize}[leftmargin=1.0em, topsep=0pt, itemsep=0pt]
    \item Optimized the classical Kalman filter algorithm in the field of motor control for brain-computer interfaces, improving computational efficiency, accuracy, and interpretability in motor decoding.
    \item Implemented high-temporal-resolution motor decoding using spiking neural networks.
\end{itemize}

\textbf{Non-invasive electroencephalogram signal decoding algorithms designing} \hfill 2021-2024
\begin{itemize}[leftmargin=1.0em, topsep=0pt, itemsep=0pt]
    \item Combined deep learning models such as graph neural networks and convolutional neural networks to achieve EEG decoding for applications including motor imagery and epilepsy detection.
    \item Investigated feature extraction to derive high-signal-to-noise-ratio features from EEG signals.
\end{itemize}


\section*{Technical Skills}

\textbf{Programming:} Python, MATLAB

\textbf{Languages:} Chinese (Native), English (IELTS 7.0), Cantonese (Beginner)


\end{document}
